% ============================================================================
% kapitel/axiome.tex
% ============================================================================

\section{Die Axiome der IWT}

Die IWT hat neun Axiome.

\subsection{Axiom 1: Diskretes Informationsfeld}

Es existiert ein skalares, diskretes Informationsfeld \( I_k^{(n)} \).

\subsection{Axiom 2: Informationserhaltung}

\[
\sum_{k=1}^{N} |I_k^{(n)}|^2 = \text{const}
\]

\subsection{Axiom 3: Informationsmetrik}

Die Metrik emergiert aus der Kopplungsmatrix:

\[
g_{kl}^{(n)} = \frac{K_{kl}^{(n)}}{\sqrt{K_{kk}^{(n)} K_{ll}^{(n)}}}
\]

\subsection{Axiom 4: Variationsprinzip}

Die Dynamik folgt aus einem diskreten Lagrange-Funktional.

\subsection{Axiom 5: Dynamische Metrik}

Die Metrik entwickelt sich dynamisch.

\subsection{Axiom 6: Energie als emergenter Informationsfluss}

Energie ist eine Erhaltungsgröße des Informationsflusses.

\subsection{Axiom 7: Naturkonstanten als emergente Skalierungsparameter}

Die Naturkonstanten folgen aus der Informationsdynamik.

\subsection{Axiom 8: Emergenz von Raum, Zeit und Dynamik}

Raum, Zeit und Dynamik emergieren aus der Information.

\subsection{Axiom 9: Universelle Gültigkeit}

Die IWT gilt auf allen Skalen.

\section{Q ist kein Axiom}

Q ist nicht in den Axiomen enthalten.

Q ist ein Resultat der Selbstorganisation.

Q ist die Quelle.

Q emergiert aus der Information.